This chapter presents the design of each subsystem of the voice-controlled drone system.
As shown in Figure~\ref{fig:block-diagram}, the system comprises three subsystems: the voice
control unit, which converts spoken commands into CRSF RC frames; the ExpressLRS (ELRS) RF
communication link, which carries those frames wirelessly to the aircraft; and the drone flight
system, which executes the commands.
The following sections describe the design procedure and implementation details for each subsystem.

%------------------------------------------------------------------------------
\subsection{Voice Control Subsystem}
\label{subsec:voice}
%------------------------------------------------------------------------------

\subsubsection{Design Procedure}

The ground-side control subsystem must convert operator intent into RC channel values
compatible with the CRSF protocol and transmit them to the airborne flight controller at a
sufficient rate to maintain a stable RF link.
Two interaction paradigms were evaluated: continuous gesture-based input derived from inertial
sensors, and discrete voice command recognition.

Gesture-based input with an IMU offers proportional control authority but requires the operator
to maintain deliberate wrist posture throughout flight, increasing physical fatigue and the risk of
inadvertent inputs.
Voice command recognition maps a fixed vocabulary of spoken words to pre-defined flight
maneuvers, reducing the physical interface to speech and freeing the operator's hands.
For a proof-of-concept demonstration in which the primary goal is autonomous takeoff,
directional movement, and landing, a discrete command set is sufficient and voice recognition
was adopted.

\subsubsection{Hardware}

The control unit is built on an ESP32 microcontroller module, selected for its dual hardware
UART peripherals, 240~MHz Xtensa LX6 core, and broad Arduino ecosystem support.
A VC02 offline voice recognition module is connected to the ESP32 on UART2 (ESP32 GPIO~16
RX, GPIO~17 TX) at 115,200~bps.
The VC02 performs on-device keyword spotting without a cloud connection, returning a single
command byte over UART each time a recognized phrase is detected.
A 128$\times$64 pixel OLED display (SSD1306, I\textsuperscript{2}C address 0x3C) is connected
on the I\textsuperscript{2}C bus and shows the current action state, barometric altitude, and CRSF
link status to the ground operator in real time.

\subsubsection{Voice Command Set and Channel Mapping}

The VC02 is configured to recognize seven command words.
Table~\ref{tab:voice-cmds} lists each command, the byte value returned by the VC02, and the
resulting RC channel action produced by the firmware.

\begin{table}[htb!]
    \centering
    \caption{Voice command set and corresponding flight actions.}
    \label{tab:voice-cmds}
    \begin{tabularx}{\linewidth}{c l X}
        \hline
        \textbf{VC02 byte} & \textbf{Command} & \textbf{Flight action} \\
        \hline
        1 & Takeoff & Arm motors, ramp throttle: 1250~\textmu s for 1.5~s, then 1180~\textmu s
            for 1.5~s, settle at hover throttle 1100~\textmu s \\
        2 & Land    & Reduce throttle gradually to 1050~\textmu s over 2~s, then disarm \\
        3 & Forward & Pitch channel to 1550~\textmu s (forward tilt) for 1200~ms, return to
            neutral \\
        4 & Backward & Pitch channel to 1450~\textmu s (backward tilt) for 1200~ms, return to
            neutral \\
        5 & Left    & Yaw channel held at 1450~\textmu s until next command \\
        6 & Right   & Yaw channel held at 1550~\textmu s until next command \\
        7 & Stop / Brake & If moving forward: reverse pitch pulse for braking; if moving
            backward: forward pitch pulse; otherwise: hold hover \\
        \hline
    \end{tabularx}
\end{table}

All channel values are in standard RC pulse-width units (\textmu s).
The neutral (mid-stick) value is 1500~\textmu s and the arm/disarm channel (CH5) is toggled
between 1000~\textmu s (disarmed) and 2000~\textmu s (armed).
Commands 3--7 are accepted only when the aircraft is in the flying state; any such command
issued before takeoff is silently discarded to prevent unintended ground arming.

\subsubsection{CRSF Frame Generation and Telemetry}
\label{subsubsec:crsf-fw}

The ESP32 encodes 16 RC channels into CRSF frames and transmits them to the RadioMaster RP2
on UART1 (ESP32 GPIO~25 TX, GPIO~26 RX) at 115,200~bps.
Each 11-bit channel value is linearly mapped from the 1000--2000~\textmu s range to the CRSF
integer range 192--1791 using the relation

\begin{equation}
    v_{\text{CRSF}} = \frac{(v_{\mu s} - 1000) \times 1600}{1000} + 192
\end{equation}

The 16 values are then packed into 22 bytes with 11 bits per channel (176 bits total), appended
with a DVB-S2 CRC-8 byte computed over the frame type and payload fields, and wrapped in the
26-byte CRSF envelope (address 0xC8, length byte, type 0x16).
Frames are sent every 20~ms (50~Hz), the minimum rate required to prevent the ELRS receiver
from triggering a link-loss failsafe.

On the downlink, the firmware parses incoming CRSF telemetry frames using a byte-at-a-time
state machine that resets whenever a gap of more than 5~ms is detected between consecutive
bytes, providing robustness against partial frames caused by UART contention.
Barometric altitude telemetry (CRSF frame type 0x09) is decoded and displayed on the OLED
in real time, giving the operator situational awareness without a separate ground station laptop.

%------------------------------------------------------------------------------
\subsection{ELRS RF Communication Subsystem}
\label{subsec:elrs}
%------------------------------------------------------------------------------

\subsubsection{Design Procedure}

Three 2.4 GHz RC link protocols were evaluated: Futaba S-FHSS, FrSky ACCESS, and
ExpressLRS (ELRS).
Table~\ref{tab:rf-comparison} summarizes the key characteristics of each.

\begin{table}[htb!]
    \centering
    \caption{Comparison of evaluated 2.4 GHz RC link protocols.}
    \label{tab:rf-comparison}
    \begin{tabularx}{\linewidth}{l c c c c}
        \hline
        \textbf{Protocol} & \textbf{Max update rate} & \textbf{Min latency} &
            \textbf{Open source} & \textbf{TX power} \\
        \hline
        Futaba S-FHSS  & 55 Hz    & $\approx$18 ms & No  & Fixed      \\
        FrSky ACCESS   & 150 Hz   & $\approx$7 ms  & No  & Fixed      \\
        ELRS 2.4 GHz   & 1000 Hz  & $<$5 ms        & Yes & Configurable \\
        \hline
    \end{tabularx}
\end{table}

Futaba S-FHSS operates at a fixed 55 Hz and uses a closed firmware ecosystem, precluding
integration with custom transmitter hardware.
FrSky ACCESS supports up to 150 Hz but similarly relies on proprietary firmware.
ELRS is an open-source protocol supporting update rates from 50 Hz to 1000 Hz with a publicly
documented firmware stack that can be compiled and deployed on commodity hardware.
ELRS was selected for its low latency, configurable update rate, and compatibility with custom
transmitter hardware.

For the transmitter unit, two hardware options were considered.
The first is a dedicated ELRS TX module (e.g., RadioMaster Ranger Micro), which provides a
self-contained solution with a preconfigured USB interface.
The second is a RadioMaster RP2 receiver module reflashed with ELRS transmitter firmware.
A dedicated TX module adds unnecessary cost and physical bulk for a fixed ground station
installation.
The RP2 module, when flashed with the \emph{RX as TX} firmware variant, performs identically as a
transmitter while occupying a smaller footprint at lower cost.
The RP2-as-TX approach was adopted.

\subsubsection{Hardware}

The RadioMaster RP2 module contains two primary integrated circuits.
The Espressif ESP8285 is a single-chip microcontroller integrating a 32-bit Tensilica L106 core
and a 2.4 GHz 802.11b/g/n Wi-Fi radio.
In the ELRS firmware context, the ESP8285 executes the protocol stack, manages the CRSF
serial interface to the host device, and controls the RF transceiver over SPI.
The Semtech SX1280 is a 2.4 GHz multi-mode transceiver supporting FLRC (Fast Long-Range
Communication) modulation, which ELRS employs at high update rates to achieve packet air times
below 1 ms and end-to-end round-trip latencies below 5 ms at 250 Hz.

\subsubsection{Firmware Configuration}
\label{subsubsec:fw-config}

Firmware was compiled using the ExpressLRS Configurator application.
Three parameters were fixed at compile time.

\textbf{Firmware target.}
The target
\texttt{Generic 2.4\,GHz /\allowbreak Generic ESP8285 2.4\,GHz RX as TX} was selected.
This variant configures the ELRS firmware to initiate and maintain the RF link, assigning the
module the transmitter role in the ELRS packet timing scheme.
A standard receiver build (\texttt{Generic ESP8285 2.4\,GHz RX}) causes the module to listen
for an external transmitter signal and never initiate packets, making it unsuitable for ground-side
use.

\textbf{Regulatory domain.}
The \texttt{ISM\_2400} domain was selected in preference to \texttt{CE\_LBT}
(Listen-Before-Talk).
The CE\_LBT domain mandates that the transmitter sense channel occupancy before each
transmission and defer if the channel is busy.
This listen phase adds a variable pre-transmission delay and reduces the time available for
telemetry data within each packet interval.
At high update rates (500 Hz and 1000 Hz), the cumulative LBT overhead measurably reduces
telemetry throughput.
The ISM\_2400 domain removes this restriction, restoring full packet capacity for combined
command and telemetry data.

\textbf{Binding Phrase.}
A user-defined alphanumeric string is compiled into both the transmitter and receiver firmware
builds.
At power-on, any two ELRS devices sharing an identical Binding Phrase and a compatible major
firmware version synchronize automatically, without a physical bind-button pairing sequence.
This eliminates the multi-step binding procedure required by traditional RC protocols and reduces
the risk of inadvertent binding to a third-party transmitter.

\subsubsection{Initial Firmware Flashing}
\label{subsubsec:flashing}

The ESP8285 is programmed via UART using the flashing tool integrated into the ExpressLRS
Configurator.
The chip enters its UART bootloader only when the GPIO0 strapping pin is sampled low at
power-on.
The prescribed procedure is: (1) hold the RP2's Boot button, asserting GPIO0 low through the
onboard pull-down circuit; (2) apply power to the module; (3) release the button.
The configurator then transfers the firmware image over UART at 921,600 baud.

A hardware behavior present on some RP2 units causes the module to enter bootloader mode
automatically when a USB-to-UART adapter is connected before power is applied.
During the pre-enumeration interval, certain adapters drive their TX and RX lines low rather than
holding the idle-mark (logic high) state.
The adapter may hold its TX line low during the reset window.
That state is read as a bootloader entry request on GPIO0.
The recommended mitigation is to supply power to the RP2 first, allow it to complete its normal
boot sequence, then connect the adapter's TX and RX lines.
This ensures the strapping pins are sampled in the normal-boot state before the adapter's lines
are driven.

After the initial UART flash, subsequent firmware updates can be applied over Wi-Fi through the
module's built-in web interface (WebUI), eliminating the need for the Boot button procedure on
future updates.

\subsubsection{CRSF Interface Baud Rate}
\label{subsubsec:crsf}

The Crossfire Serial (CRSF) protocol carries RC channel data between the RP2 and the connected
host ESP32 over a full-duplex UART interface.
The baud rate must match on both sides.
A rate of 115,200~bps was used, which is the default CRSF baud rate supported by the RP2
hardware and the standard Arduino \texttt{HardwareSerial} peripheral on the ESP32.

A 26-byte CRSF frame at 115,200~bps with 8-N-1 framing occupies approximately 2.26~ms on
the wire.
Because frames are sent every 20~ms (50~Hz command rate), the duty cycle is approximately
11\%, leaving ample headroom for returning telemetry frames in the same UART window without
frame collision.

\subsubsection{Telemetry Configuration}
\label{subsubsec:telemetry}

ELRS supports bidirectional telemetry: the flight controller returns sensor data
(battery voltage, received signal strength, link quality) through reserved downlink slots.
Two operating modes were configured.

\textbf{RC control mode.}
When the RP2 transmits as the control head connected to the ESP32 ground unit, the telemetry ratio is
set via the \texttt{elrsV3.lua} script on the transmitter.
A ratio of 1:64 was configured, allocating one downlink telemetry slot for every 64 uplink RC
command packets.
At a 250 Hz command update rate this yields approximately 3.9 telemetry packets per second,
sufficient for battery voltage and link quality monitoring without perceptibly reducing command
throughput.

\textbf{Ground station serial mode.}
When the RP2 is connected by USB to a laptop running ground station software, the
\emph{Use as AirPort Serial device} option is enabled in the WebUI.
In this mode, all received CRSF telemetry frames are forwarded as raw serial data through the
module's USB-to-UART bridge at 400,000 bps, allowing ground station applications to parse and
display real-time flight data without a separate RF modem.

%------------------------------------------------------------------------------
\subsection{Drone Flight Subsystem}
\label{subsec:drone}
%------------------------------------------------------------------------------

\subsubsection{Airframe and Propulsion}

The airframe is a five-inch class quadrotor based on a Mark4-style five-inch frame
(purchased as a complete frame kit).
Four brushless outrunner motors are installed; each is rated for 1860 KV and is paired with a
blade compatible with the five-inch propeller class.
A five-inch propeller protection set was added (vendor listing: four-piece guard set for five-inch
class props) to reduce damage during low-altitude development flights.

\subsubsection{Flight Control Electronics}

The autopilot uses a SpeedyBee F405 V5 series flight controller.
Motor commutation is driven by a dedicated OX32 55 A, four-in-one ESC wired to the four
1860~KV outrunners.
During procurement the team received \emph{two} four-in-one ESCs from different manufacturers:
one shipped with the initial SpeedyBee-oriented stack bundle and the OX32 unit purchased as a
separate line item.
Because the brands differed and the newer OX32 hardware matched the final integration plan, the
earlier ESC was taken out of the flight harness, retained as laboratory spare stock, and processed for
reimbursement under the spare-equipment policy, while the OX32 module remained in the
demonstrator.

An ExpressLRS compatible airborne receiver (RX24T) terminates the RF link from the ground-side
RadioMaster RP hardware described in Section~\ref{subsec:elrs}.
The flight controller communicates with that receiver over the same CRSF serial interface
convention used throughout the ELRS ecosystem.

\paragraph{Procurement note.}
Vendor invoices occasionally abbreviate bundled boards with non-specific stock codes while still
shipping SpeedyBee-class hardware; the installed configuration documented here was verified on
the bench and in flight, not inferred from invoice text alone.

\subsubsection{Commissioning and Calibration}

Before first flight the following commissioning steps were performed using ground station software
connected to the SpeedyBee F405 V5 over USB.

\textbf{Sensor calibration.}
Accelerometer calibration was carried out by placing the airframe on a flat surface and following
the six-position guided routine in the ground station software; each face of the airframe was held
still while the software recorded the gravity vector.
Gyroscope calibration was performed with the airframe stationary on a bench to establish the
zero-rate offset for each axis.
The flight controller stores both calibration sets in onboard flash and reloads them at power-on.

\textbf{Motor direction verification.}
Each motor was spun individually at minimum throttle through the ground station motor test page.
Two motors were found to rotate in the incorrect direction for the chosen propeller-rotation scheme
and were reversed by swapping any two of the three ESC phase wires at the ESC pads.
After swapping, all four motors were confirmed to spin in the intended direction before propellers
were fitted.

\textbf{Flight controller configuration.}
The SpeedyBee F405 V5 ships with a Betaflight firmware image.
The factory PID values were retained for all three attitude axes; no manual PID tuning loop was
performed.
This decision was informed by the fact that a five-inch class quadrotor at 1860~KV is a
well-characterized platform, and the factory Betaflight defaults produced stable, flyable behavior
in early hover tests without further adjustment.

\subsubsection{Manual Flight Qualification}

A RadioMaster Pocket transmitter was used for manual flight qualification prior to integrating the
smartwatch control path.
The pilot adapted to the aircraft's handling characteristics over multiple hover and forward-flight
sessions, verifying stable attitude hold in angle mode.
This baseline established a reference for evaluating whether the smartwatch control path
introduced any perceivable latency or authority degradation relative to conventional stick input.
